\documentclass[a4paper,10pt]{report}\input{../0.Préambule/Préambule_cours_prof}\begin{document}

\pagestyle{empty}
\newpage
\subsection*{Activité : Racine carrée d'un nombre}
\addcontentsline{toc}{subsection}{Activité : Racine carrée d'un nombre}

\begin{enumerate}[font=\color{exer1}]
\item Comment se lit \og  $7$m$^2$ \fg{} ?

\medskip
\dotfill
\item Dans un carré $ABCD$, de côté $5$cm, on souhaite déterminer son aire.
\begin{enumerate}[font=\color{exer1}]
\item Quelle est la formule de l'aire d'un carré de côté $c$ ? 

\medskip
\dotfill
\item Quelle est l'aire d'un carré de côté $5$cm ?

\medskip
\dotfill
\item Compléter les tableaux suivants.

\renewcommand{\arraystretch}{1.8}
\begin{tabularx}{\linewidth}{|>{\columncolor{gray!20}}p{4cm}|*{10}{>{\centering \arraybackslash}X|}}
\hline
Côté d'un carré (cm) & $2$ & $9$ & $2,5$ & $3,2$ & $13$ & \dotfill & \dotfill & \dotfill & \dotfill & \dotfill \\\hline
Aire du carré (cm$^2$) & \dotfill & \dotfill & \dotfill & \dotfill & \dotfill & $1$ & $9$ & $25$ & $36$ & $100$\\\hline
\end{tabularx}
\end{enumerate}
\item On souhaite déterminer le côté d'un carré d'aire $20$cm$^2$.
\begin{enumerate}[font=\color{exer1}]
\item Déterminer un encadrement par des valeurs entières du côté du carré.

\medskip
\dotfill
\item à l'aide de la touche $\boxed{\displaystyle\sqrt{\quad}}$ de la calculatrice, déterminer $\displaystyle\sqrt{20}$ puis écrire le résultat en arrondissant au centième.

\medskip
\dotfill
\medskip
\item On peut répondre à la question initiale, le côté d'un carré d'aire $20$cm$^2$ est de \dotfill
\item Compléter le tableau suivant :
\renewcommand{\arraystretch}{1.8}
\begin{tabularx}{10.7cm}{|>{\columncolor{gray!20}}p{1cm}|*{5}{>{\centering \arraybackslash}X|}}
\hline
$c^2$ & $20$ & $2$ & $6$ & $5$ & $49$\\\hline
$c$ & \dotfill & \dotfill & \dotfill & \dotfill & \dotfill\\\hline
\end{tabularx}
\end{enumerate}
\end{enumerate}

\end{document}